% =============================================================
% CAPÍTULO 1 — Introducción
% Versión: matrícula de honor
% =============================================================

\section{Introducción}
\label{sec:introduccion}

\subsection{Contexto y motivación}
\label{sec:contexto-motivacion}

Los modelos de lenguaje basados en la arquitectura Transformer
\cite{vaswani2017attention} han transformado el procesamiento del lenguaje natural
en la última década. BERT \cite{devlin2019bert}, con sus 110~millones de parámetros,
demostró en 2018 que una sola red pre-entrenada podía alcanzar resultados de
referencia en tareas tan diversas como clasificación de texto, extracción de
información o respuesta a preguntas. Hoy, 110~millones de parámetros es la talla
\emph{pequeña}: GPT-3 alcanzó los 175\,000~millones \cite{brown2020language}, y los
modelos abiertos recientes operan en el rango de 8\,000 a 405\,000~millones. El coste
de desplegarlos en dispositivos móviles, sistemas embebidos o infraestructuras con
restricciones de latencia y energía sigue siendo un problema abierto, y proyecciones
recientes estiman que la inferencia ---no el entrenamiento--- concentra el grueso del
gasto energético a largo plazo de un modelo desplegado \cite{strubell2019energy,
schwartz2020green}.

Frente a este escenario, una pregunta directa orienta este trabajo: ¿se puede hacer
un modelo más pequeño sin destruir lo que ha aprendido? La compresión de modelos
lleva años intentando responderla con cuatro herramientas principales ---poda,
cuantización, destilación y factorización de bajo rango. Aquí se utiliza la última
en su variante más clásica: la descomposición en valores singulares (SVD), aplicada
de forma selectiva a las 72 matrices lineales del encoder de BERT-base. La SVD
ofrece tres ventajas para los objetivos de este trabajo: no requiere
re-entrenamiento, admite granularidad arbitraria (rango específico por matriz), y
proporciona garantías matemáticas de optimalidad respecto a la norma de Frobenius
mediante el teorema de Eckart-Young \cite{eckart1936approximation}.

Pero el proyecto no se limita a la compresión. Lo que comenzó como un análisis de
cuánto puede comprimirse cada componente derivó en algo más interesante: al observar
qué se rompe y qué sobrevive a la eliminación selectiva de partes del modelo, se
obtiene un mapa empírico de cómo está organizada internamente la información. La
compresión funciona, en este sentido, como un \emph{bisturí experimental}: cada
configuración comprimida es una intervención controlada que revela la importancia
funcional del componente afectado. Esta perspectiva conecta directamente la
compresión con la \emph{interpretabilidad mecánica}, una línea de investigación
reciente que busca abrir el modelo y caracterizar los circuitos internos
responsables de comportamientos específicos \cite{elhage2021mathematical,
meng2022locating, conmy2023towards}.

El trabajo combina así dos perspectivas habitualmente disjuntas en la literatura.
Por un lado, la compresión SVD selectiva por componente y profundidad. Por otro,
cinco técnicas de interpretabilidad mecánica ---probing lineal, \emph{logit lens},
\emph{activation patching}, ablación de cabezas de atención y análisis de selectividad
neuronal en las redes feed-forward--- que examinan la organización interna del
modelo desde ángulos diferentes pero convergentes. Y cierra el bucle traduciendo
esa comprensión en una estrategia de compresión \emph{informada por datos} que
domina la frontera de Pareto frente a las aproximaciones ciegas.

La elección de la clasificación de emociones como dominio de aplicación no es
arbitraria. GoEmotions \cite{demszky2020goemotions} ofrece 23~categorías
emocionales con propiedades semánticas heterogéneas: algunas emociones tienen
vocabulario distintivo (\emph{gratitude} $\leftrightarrow$ ``thanks''), mientras que
otras requieren comprensión pragmática profunda (\emph{disappointment},
\emph{realization}). Esta diversidad convierte la tarea en un test exigente para la
compresión, porque permite observar cómo la degradación del modelo afecta de forma
\emph{diferencial} a emociones con distintos perfiles lingüísticos. Un dataset
binario (sentimiento positivo/negativo) no haría visible esta vulnerabilidad
selectiva, que constituye uno de los hallazgos centrales del proyecto.

La motivación tiene tres vertientes. Desde el punto de vista \emph{práctico}, la
compresión SVD se aplica después del entrenamiento, no requiere hardware
especializado y es directamente escalable a cualquier modelo con capas lineales.
Desde el punto de vista \emph{científico}, la compresión selectiva proporciona
acceso a la organización interna del modelo de un modo que las técnicas
correlacionales (como el probing) no pueden replicar: cuando comprimir un componente
destruye una capacidad, esa relación es funcional, no solo asociativa. Y desde el
punto de vista de la \emph{sostenibilidad computacional}, reducir el tamaño de los
modelos para inferencia contribuye directamente a la filosofía \textit{Green AI}
\cite{schwartz2020green}.


\subsection{Formulación del problema}
\label{sec:formulacion-problema}

El problema central de este trabajo se articula en torno a tres preguntas
encadenadas. La primera describe; la segunda explica; la tercera prescribe.

\begin{quote}
    \textbf{P1 (descriptiva).} \textit{¿Cómo afecta la compresión SVD selectiva,
    aplicada por tipo de componente y por profundidad de capa, al rendimiento y a
    la calidad de las representaciones internas de un modelo Transformer
    \emph{fine-tuneado} para clasificación de emociones?}
\end{quote}

\begin{quote}
    \textbf{P2 (explicativa).} \textit{¿Qué revelan las técnicas de interpretabilidad
    mecánica sobre la organización interna de la información emocional en el
    modelo, y qué propiedades estructurales explican la sensibilidad diferencial
    observada en P1?}
\end{quote}

\begin{quote}
    \textbf{P3 (prescriptiva).} \textit{¿Puede el conocimiento mecanicista
    obtenido en P2 traducirse en estrategias de compresión que dominen a las
    aproximaciones ciegas establecidas en la literatura?}
\end{quote}

Los trabajos previos de compresión SVD en Transformers \cite{noach2020compressing,
hsu2022language} se han centrado mayoritariamente en compresión uniforme: el mismo
grado de reducción para todas las capas y componentes. Esta uniformidad ignora dos
asimetrías documentadas. Primero, las distintas capas de BERT codifican tipos de
información diferentes ---las tempranas capturan rasgos léxicos y sintácticos, las
tardías capturan rasgos semánticos y específicos de la tarea \cite{tenney2019bert,
jawahar2019what}. Segundo, los componentes de atención y las redes feed-forward
desempeñan roles funcionales distintos \cite{geva2021transformer, voita2019analyzing}.
La hipótesis subyacente del proyecto es que tratarlos igual es subóptimo, y que
medir cuánto subóptimo es ---y por qué--- requiere combinar compresión sistemática
con interpretabilidad mecánica.


\subsection{Objetivos}
\label{sec:objetivos}

El objetivo general del proyecto es:

\begin{quote}
    Producir un mapa funcional de cómo BERT-base \emph{fine-tuneado}
    organiza internamente la información emocional, utilizando la
    compresión SVD selectiva como instrumento experimental
    ---\emph{bisturí} que mide la contribución funcional de cada
    componente por la magnitud de lo que se rompe al eliminarlo--- y
    triangulando los hallazgos con cinco técnicas de interpretabilidad
    mecánica. Como aplicación derivada del mapa, se diseña una
    estrategia de compresión informada por datos empíricos de
    sensibilidad y se compara con las aproximaciones ciegas
    establecidas en la literatura.
\end{quote}

Este objetivo se descompone en cinco objetivos específicos que estructuran la
contribución y delimitan el alcance:

\begin{enumerate}[nosep, label=\textbf{O\arabic*.}]
    \item \textbf{Caracterizar la sensibilidad a la compresión SVD} por tipo de
    componente (Q, K, V, Output, FFN Intermediate, FFN Output), por profundidad
    de capa y por emoción, estableciendo un mapa cuantitativo del trade-off
    compresión--rendimiento bajo compresión uniforme y adaptativa.

    \item \textbf{Localizar y mapear la organización interna de la información
    emocional} mediante cinco técnicas de interpretabilidad mecánica
    complementarias: probing lineal para identificar la profundidad de
    cristalización de cada emoción, \emph{logit lens} para distinguir información
    presente de información alineada con la base del clasificador, activation
    patching para establecer dependencias funcionales, ablación de cabezas de
    atención para mapear la especialización atencional y análisis de selectividad
    neuronal para caracterizar la codificación en las redes feed-forward.

    \item \textbf{Establecer conexiones cuantitativas entre estructura interna y
    vulnerabilidad a la compresión}, identificando qué propiedades del modelo
    (profundidad de cristalización, intensidad de codificación neuronal,
    distribución espectral) predicen la caída de rendimiento bajo SVD.

    \item \textbf{Diseñar y evaluar una estrategia de compresión informada} que
    traduzca los hallazgos de interpretabilidad en un algoritmo de asignación de
    rangos, y demostrar su superioridad sobre las aproximaciones ciegas en la
    frontera de Pareto.

    \item \textbf{Evaluar la recuperación por fine-tuning} del modelo
    comprimido, investigando si la compresión SVD seguida de un breve
    re-entrenamiento puede actuar como regularización implícita que beneficie
    a las emociones infrarrepresentadas.
\end{enumerate}


\subsection{Contribuciones}
\label{sec:contribuciones-intro}

Las contribuciones principales de este trabajo son las siguientes:

\begin{enumerate}[label=\textbf{C\arabic*.}]

    \item \textbf{Mapa de sensibilidad a la compresión} desagregado por tipo de
    componente, profundidad de capa y emoción (O1). El mapa documenta una
    asimetría radical entre componentes: las proyecciones Query y Key toleran
    compresión a rango~128 con retención del 99\% del F1 macro, mientras que la
    FFN Intermediate al mismo rango colapsa al 7\%. La diferencia de
    sensibilidad alcanza un factor 14$\times$ en retención absoluta de F1 a
    rango~128, que se amplifica hasta un factor 72$\times$ a rango~64
    cuando se normaliza por el porcentaje de parámetros eliminados (Query
    retiene 98{,}6\%, FFN Intermediate colapsa al 0\%). Hasta donde sabemos,
    esta asimetría no había sido documentada cuantitativamente en la
    literatura sobre SVD en Transformers.

    \item \textbf{Triangulación de interpretabilidad mecánica} mediante cinco
    técnicas independientes (probing, \emph{logit lens}, \emph{activation patching}, ablación
    de cabezas y análisis neuronal) (O2). Los resultados convergen en una
    arquitectura funcional de tres niveles ---señal léxica en capas tempranas,
    cómputo de transición en capas medias, alineación geométrica con la base
    del clasificador en capas tardías--- que explica de forma unificada
    fenómenos previamente disjuntos: el patrón en U del \emph{logit lens}, la
    suficiencia causal de la capa~11 en \emph{activation patching}, el gradiente de
    criticidad atencional documentado por ablación, y la concentración del 84\%
    de las neuronas significativas en las capas 8--11.

    \item \textbf{Resolución de la paradoja información-dependencia.} El
    probing demuestra que las capas tardías apenas aportan información nueva
    (ganancia $<0{,}004$ en L10--L11), pero el \emph{activation patching} las
    identifica como las únicas suficientes para restaurar el modelo desde un
    colapso total. Mediante el \emph{logit lens} documentamos que esta aparente
    contradicción refleja una distinción funcional precisa: las capas tardías
    no \emph{crean} señal emocional ---la \emph{rotan} hacia la base sobre la
    que opera la cabeza de clasificación. La distinción entre información
    \emph{presente} (probing), información \emph{leíble} (\emph{logit lens}) y
    componentes \emph{suficientes} (patching) constituye una clarificación
    metodológica que va más allá del caso estudiado.

    \item \textbf{Conexión cuantitativa entre interpretabilidad y compresión}
    (O3). La profundidad de cristalización por probing predice la
    vulnerabilidad a la compresión SVD a rango 128
    (Spearman~$\rho \approx 0{,}48$), y la norma del vector de selectividad
    neuronal es el mejor predictor individual de la caída de F1 bajo
    compresión ($\rho = 0{,}64$, $p = 0{,}001$, $n=23$). La
    geometría que la SVD optimiza es \emph{ortogonal} a la organización
    funcional del modelo (correlación entre fracción de señal eliminada y
    caída de F1: $\rho \approx 0$): la SVD no daña selectivamente las neuronas
    emocionales, sino que las daña uniformemente, y la vulnerabilidad
    diferencial proviene de la intensidad de codificación, no de la geometría
    SVD.

    \item \textbf{Algoritmo greedy de compresión informada por datos} (O4),
    alimentado con datos empíricos de sensibilidad obtenidos en C1, que domina
    la frontera de Pareto con 8 de los 9 puntos Pareto-óptimos del estudio. A
    ratio de parámetros del 80\%, la estrategia greedy retiene el 87\% del F1
    macro, mientras que la compresión uniforme al mismo ratio retiene solo el
    43\%. El algoritmo descubre de forma independiente la jerarquía de
    importancia identificada por interpretabilidad ---comprime primero Q y K,
    nunca toca la FFN Intermediate de las capas tardías---, lo que constituye
    una validación cruzada entre dos líneas metodológicamente independientes.

    \item \textbf{Resultado negativo informativo: la heurística cualitativa
    no extrae valor adicional.} Tres estrategias informadas basadas en reglas
    heurísticas derivadas de los hallazgos de interpretabilidad
    \emph{convergen exactamente} sobre la frontera de las baselines ciegas
    (\texttt{informed\_aggressive} reproduce \texttt{uniform\_r256},
    \texttt{informed\_moderate} reproduce \texttt{uniform\_r512}), o requieren
    más parámetros que el modelo original sin alcanzar las retenciones más
    altas (\texttt{informed\_light} a ratio 1{,}285). Este resultado delimita
    el valor aplicado de la interpretabilidad mecánica para la compresión: el
    conocimiento cualitativo identifica \emph{qué} medir, pero solo los datos
    cuantitativos de sensibilidad permiten determinar \emph{cuánto}
    comprimir. La narrativa completa (heurística no innova $\rightarrow$
    pivot data-driven $\rightarrow$ dominancia Pareto) es, en sí misma, una
    observación metodológica transferible.

    \item \textbf{Evidencia preliminar de regularización por compresión}
    (O5). Un modelo comprimido al 86\% de parámetros, tras tres épocas de
    \emph{fine-tuning}, alcanza un F1 macro de 0{,}591 frente a 0{,}577 del baseline
    sin comprimir. Las ganancias se concentran en emociones
    infrarrepresentadas: \emph{embarrassment} pasa de F1 = 0{,}267 a F1 =
    0{,}509 ($+90\%$ relativo). Reportamos este hallazgo como observación
    consistente con la hipótesis de que la SVD, al eliminar direcciones de
    baja energía, actúa como regularización estructural que reduce el
    sobreajuste a patrones espurios. El alcance experimental (un modelo, una
    tarea, sin grupo de control con epochs adicionales) impide afirmar
    causalidad: en el Capítulo~\ref{sec:conclusiones} se discuten los
    experimentos de control que verificarían la hipótesis.

\end{enumerate}


\subsection{Marco del proyecto en la FIB}
\label{sec:marco-fib}

Este proyecto se enmarca en el Grado en Ingeniería Informática de la Facultat
d'Informàtica de Barcelona (FIB-UPC), dentro de la especialidad de Computación.
El trabajo integra conocimientos de múltiples asignaturas del grado. La
descomposición SVD y el análisis espectral se fundamentan en álgebra lineal y
cálculo matricial (M1, M2), mientras que la estabilidad numérica de las
aproximaciones de bajo rango conecta directamente con los contenidos de
Computación Numérica (CN). El diseño y la evaluación del algoritmo greedy de
compresión aplican técnicas de diseño algorítmico y análisis de complejidad (A,
EDA). Los fundamentos de aprendizaje automático (IA) sustentan tanto el
\emph{fine-tuning} de BERT como los clasificadores de probing, y la ejecución
eficiente en GPU (PAR) ha sido necesaria para el entrenamiento y la evaluación
sobre el dataset completo. Las metodologías de desarrollo de software (PROP)
se reflejan en la organización modular del código experimental. Un intercambio
en la Technical University of Denmark (DTU, Q1 2024--25) proporcionó formación
complementaria de nivel máster en \textit{Machine Learning}, \textit{Deep
Learning} y \textit{Graph Neural Networks}, incluyendo contacto con técnicas de
interpretabilidad mecánica que han sido instrumentales en el diseño del
Capítulo~\ref{sec:resultados-interpretabilidad}.

El proyecto tiene carácter eminentemente experimental y de investigación: no
se desarrolla un producto de software, sino que se lleva a cabo un estudio
sistemático sobre cómo la compresión SVD afecta a los distintos componentes de
un Transformer y qué revela esto sobre la organización interna de las
representaciones emocionales. La metodología es comparable a la de la
investigación empírica en NLP: hipótesis explícitas, experimentos controlados,
métricas estándar (F1 macro y micro, retención, ratio de parámetros), análisis
estadístico (correlaciones de Spearman con reporte de p-valores) y reporte
honesto de resultados ---incluidos los nulos.


\subsection{Alcance y limitaciones explícitas}
\label{sec:alcance-limitaciones}

Por transparencia metodológica, los resultados del trabajo deben interpretarse
dentro de las siguientes restricciones, que se discuten en mayor detalle en el
Capítulo~\ref{sec:conclusiones}:

\begin{itemize}[nosep]
    \item \textbf{Modelo único}: los experimentos se realizan sobre BERT-base
    (110\,M parámetros). La generalización a BERT-large, RoBERTa, o
    arquitecturas decoder-only (GPT-2, LLaMA) no se verifica empíricamente
    en este trabajo.

    \item \textbf{Tarea única}: clasificación multi-etiqueta de emociones
    sobre GoEmotions. La extensión a NER, inferencia textual o tareas
    generativas queda como trabajo futuro.

    \item \textbf{Compresión post-hoc}: la SVD se aplica una vez fine-tuneado
    el modelo, sin compresión durante el entrenamiento ni cuantización
    posterior.

    \item \textbf{Potencia estadística limitada}: las correlaciones se
    calculan sobre $n = 23$ emociones, lo que permite detectar efectos
    grandes con confianza pero no efectos moderados.
\end{itemize}

Estas limitaciones se asumen por diseño: el alcance del trabajo prioriza la
profundidad analítica sobre un único modelo y tarea frente a estudios más
superficiales sobre múltiples configuraciones.


\subsection{Estructura de la memoria}
\label{sec:estructura-memoria}

La memoria se organiza en siete capítulos. El Capítulo~\ref{sec:estado-arte}
revisa el estado del arte en compresión de Transformers e interpretabilidad
mecánica, sitúa los fundamentos matemáticos de la SVD y posiciona el trabajo
respecto a la literatura previa. El Capítulo~\ref{sec:metodologia-tfg} describe
la metodología experimental: dataset, modelo, técnicas de compresión y
herramientas de interpretabilidad.

Los tres capítulos siguientes contienen los resultados, en orden de granularidad
creciente. El Capítulo~\ref{sec:resultados-compresion} presenta los resultados
de compresión SVD: análisis espectral, transición de fase bajo compresión
uniforme, sensibilidad por componente y profundidad, y compresión adaptativa
energética. El Capítulo~\ref{sec:resultados-interpretabilidad} expone los
resultados de las cinco técnicas de interpretabilidad mecánica y argumenta su
convergencia en una arquitectura funcional unificada. El
Capítulo~\ref{sec:sintesis} cierra el bucle: utiliza los hallazgos de los
capítulos anteriores para diseñar y evaluar la estrategia greedy de compresión
informada, documenta el resultado negativo de la heurística previa y analiza la
recuperación mediante \emph{fine-tuning}.

El Capítulo~\ref{sec:conclusiones} sintetiza las contribuciones, discute las
limitaciones y propone líneas de investigación futura. La gestión del proyecto
(planificación temporal, presupuesto, sostenibilidad) se incluye como anexo
final.